\section{System Architecture}
\label{sec:architecture}
\begin{figure*}[!t]
\centering
\begin{tikzpicture}[
  font=\scriptsize,
  box/.style={draw, rounded corners=2pt, align=center, minimum height=7mm, inner sep=3pt, fill=black!3},
  grp/.style={draw, dashed, rounded corners=4pt, inner sep=7pt},
  flow/.style={-{Latex[length=2mm]}, thick}
]
% ---- Frontend ----
\node[box] (home)  {Home / Campus};
\node[box, right=4mm of home] (tour) {Virtual Tour\\(Pannellum)};
\node[box, right=4mm of tour] (map) {Map\\(Google Satellite)};
\node[box, right=4mm of map] (chat) {Chat + Floating\\Assistant};
\node[box, right=4mm of chat] (voice) {Voice (Web\\Speech API)};
\begin{scope}[on background layer]
\node[grp, fit=(home)(tour)(map)(chat)(voice), label={[anchor=north west]north west:\textbf{Frontend --- Next.js 15 / React 19 / TypeScript}}] (fe) {};
\end{scope}

% ---- API layer ----
\node[box, below=16mm of home, xshift=6mm] (chatapi) {\code{POST /api/v1/chat}\\(sync handler)};
\node[box, right=6mm of chatapi] (tourapi) {\code{/api/v1/tour/*}\\scenes, neighbors};
\node[box, right=6mm of tourapi] (navapi) {\code{GET /api/v1/navigate}\\(no FE consumer)};
\node[box, right=6mm of navapi] (crud) {CRUD routers\\nodes, edges, \dots};
\begin{scope}[on background layer]
\node[grp, fit=(chatapi)(tourapi)(navapi)(crud), label={[anchor=north west]north west:\textbf{API --- FastAPI (Pydantic response models)}}] (api) {};
\end{scope}

% ---- Orchestration ----
\node[box, below=14mm of chatapi, xshift=-4mm] (ctx) {Context resolver\\(follow-ups)};
\node[box, right=5mm of ctx] (sup) {Supervisor\\(det.\ rules + LSTM)};
\node[box, right=5mm of sup] (agents) {5 specialist agents\\(shared pipeline)};
\begin{scope}[on background layer]
\node[grp, fit=(ctx)(sup)(agents), label={[anchor=north west]north west:\textbf{Orchestration --- \code{scripts/ai/} (imported by API)}}] (orch) {};
\end{scope}

% ---- Services ----
\node[box, below=14mm of ctx, xshift=-2mm] (retr) {Hybrid retrieval\\(dense + BM25)};
\node[box, right=4mm of retr] (rr) {Reranker\\(4-feature)};
\node[box, right=4mm of rr] (conf) {Confidence\\gate};
\node[box, right=4mm of conf] (gen) {LLM generator\\+ grounding check};
\node[box, right=4mm of gen] (nav) {A$^{\star}$ pathfinder\\+ graph builder};
\node[box, right=4mm of nav] (lstm) {LSTM intent\\classifier};
\begin{scope}[on background layer]
\node[grp, fit=(retr)(rr)(conf)(gen)(nav)(lstm), label={[anchor=north west]north west:\textbf{Domain services}}] (svc) {};
\end{scope}

% ---- Data / infra ----
\node[box, below=13mm of retr, xshift=6mm] (pg) {PostgreSQL 16\\campus schema, sessions};
\node[box, right=8mm of pg] (chroma) {ChromaDB\\\code{gat\_kb} (HNSW cosine)};
\node[box, right=8mm of chroma] (ollama) {Ollama\\Llama~3.2};
\node[box, right=8mm of ollama] (files) {Panorama /\\media assets};
\begin{scope}[on background layer]
\node[grp, fit=(pg)(chroma)(ollama)(files), label={[anchor=north west]north west:\textbf{Data \& infrastructure (Docker Compose)}}] (data) {};
\end{scope}

% edges
\draw[flow] (fe) -- node[right,font=\scriptsize]{HTTPS / axios (REST, no streaming)} (api);
\draw[flow] (api) -- (orch);
\draw[flow] (orch) -- (svc);
\draw[flow] (retr) -- (chroma);
\draw[flow] (retr.south west) to[bend right=12] (pg.north west);
\draw[flow] (gen) -- (ollama);
\draw[flow] (nav) -- (pg);
\draw[flow] (tourapi.south) to[bend left=10] (pg.north);
\end{tikzpicture}
\caption{\system{} system architecture, reconstructed from the implemented
import graph. The RAG/agent pipeline lives in \code{scripts/ai/} and is
imported into the FastAPI process via a \code{sys.path} insertion; the
navigation engine lives in \code{backend/app/navigation/}. The chat handler
is a synchronous route dispatched on Starlette's thread pool. There is
currently no front-end consumer of \code{/api/v1/navigate}.}
\label{fig:arch}
\end{figure*}

\begin{figure}[!t]
\centering
\begin{tikzpicture}[
  font=\scriptsize,
  node distance=3.2mm,
  b/.style={draw, rounded corners=2pt, align=center, text width=6.1cm, minimum height=6mm, inner sep=3pt, fill=black!3},
  d/.style={draw, diamond, aspect=2.2, align=center, inner sep=1pt, text width=2.6cm, fill=black!5},
  pf/.style={-{Latex[length=1.8mm]}}
]
\node[b] (inp) {Browser: \code{axios POST /api/v1/chat} $\{message, session\_id?\}$};
\node[b, below=of inp] (sess) {Resolve / create session (PostgreSQL, best-effort)};
\node[d, below=of sess] (amb) {ambiguous\\follow-up?};
\node[b, below=6mm of amb] (route) {\code{supervisor.route(effective\_message)}: classify $\rightarrow$ dispatch};
\node[b, below=of route] (retr) {Hybrid retrieval (top-20$\times$2) $\rightarrow$ rerank (top-5) $\rightarrow$ context selection};
\node[d, below=of retr] (cat) {confidence\\category?};
\node[b, below=8mm of cat] (llm) {Ollama \code{ChatOllama(llama3.2)} with grounding-only prompt};
\node[d, below=of llm] (grd) {numeric claims\\grounded?};
\node[b, below=8mm of grd] (out) {\code{ChatResponse}: answer, confidence, agent, sources[], nav?, panorama?};
\node[b, right=5mm of cat, text width=1.9cm, yshift=-2mm] (refuse) {LOW $\Rightarrow$ fixed refusal (LLM \emph{not} called)};
\node[b, right=5mm of grd, text width=1.9cm] (with) {withhold $\Rightarrow$ substitute refusal};
\node[b, right=5mm of amb, text width=1.9cm] (clar) {clarifying question; skip router};

\draw[pf] (inp)--(sess);
\draw[pf] (sess)--(amb);
\draw[pf] (amb)-- node[left]{no} (route);
\draw[pf] (amb)-- node[above]{yes} (clar);
\draw[pf] (route)--(retr);
\draw[pf] (retr)--(cat);
\draw[pf] (cat)-- node[left]{MED/HIGH} (llm);
\draw[pf] (cat)-- node[above]{LOW} (refuse);
\draw[pf] (llm)--(grd);
\draw[pf] (grd)-- node[left]{yes} (out);
\draw[pf] (grd)-- node[above]{no} (with);
\draw[pf] (clar) |- (out);
\draw[pf] (refuse) |- (out);
\draw[pf] (with) |- (out);
\end{tikzpicture}
\caption{End-to-end request flow for a chat query. The confidence gate and
the post-generation grounding check are two independent points at which the
system substitutes a fixed refusal for a generated answer.}
\label{fig:reqflow}
\end{figure}


\system{} is a four-tier web application (Fig.~\ref{fig:arch}): a
single-page front end, a versioned HTTP API, an orchestration layer, and a
set of domain services backed by relational, vector, model-serving, and
file infrastructure. The architecture was reconstructed for this paper
directly from the implemented import graph rather than from design
documents, several of which are stale relative to the code.

\subsection{Tier 1 --- Front End}
The front end is a Next.js~15 App-Router application in TypeScript
(\code{strict} mode). Routes are thin compositions of feature components;
all HTTP traffic is funnelled through a single \code{axios} client layer;
server state is held in TanStack Query caches and cross-component client
state in three Zustand stores (chat, campus, language). The five
user-facing surfaces are the landing/campus pages, the Pannellum-based
360$^{\circ}$ tour, a Google-satellite map view, a full-page chat, and a
site-wide floating assistant that mounts the same chat component in a
compact panel. Voice input is a browser-native Web Speech API hook that
writes its transcript into the ordinary chat send path.

\subsection{Tier 2 --- API}
The API is a FastAPI application. Every route declares a Pydantic
\code{response\_model}; cross-origin access is restricted to configured
origins. The chat route, \code{POST /api/v1/chat}, is deliberately a
\emph{synchronous} handler: the entire downstream pipeline (ChromaDB
queries, sentence-transformer encoding, SQLAlchemy calls, the Ollama HTTP
call) is blocking code, and Starlette dispatches synchronous handlers on a
thread pool, so a slow request does not block the event loop. Tour data is
served by \code{/api/v1/tour/*}; the campus graph is exposed through
generic CRUD routers; and \code{GET /api/v1/navigate} exposes the
A$^{\star}$ engine (it currently has no front-end consumer).

\subsection{Tier 3 --- Orchestration}
The RAG and multi-agent logic resides in \code{scripts/ai/} and is imported
into the FastAPI process by inserting that directory on \code{sys.path}.
This is a documented deviation from the repository's own folder convention,
made so that each pipeline stage remains independently runnable and
testable outside a live server. The chat handler performs three steps:
(i)~best-effort session resolution against PostgreSQL; (ii)~contextual
follow-up resolution, which may short-circuit the request with a clarifying
question when a pronominal reference is ambiguous; and (iii)~a single call
into \code{supervisor.route()}, which is the sole path from the HTTP layer
into the AI pipeline. Figure~\ref{fig:reqflow} shows the full flow.

\subsection{Tier 4 --- Domain Services and Infrastructure}
Six domain services are invoked by the orchestration layer: hybrid
retrieval, reranking, confidence estimation, grounded generation (with the
post-hoc grounding check), the A$^{\star}$ pathfinder with its graph
builder, and the LSTM intent classifier. Infrastructure comprises
PostgreSQL~16 (the campus schema and chat-session persistence), an
on-disk ChromaDB collection with an HNSW cosine index, an Ollama server
hosting Llama~3.2, and a panorama/media asset tree. Docker Compose defines
exactly four services (\code{frontend}, \code{backend}, \code{db},
\code{ollama}); there is no cache or queue service.

\subsection{Shared Campus Model}
All three subsystems key off one relational model:
$\textsf{Campus}\!\rightarrow\!\textsf{Building}\!\rightarrow\!\textsf{Floor}\!\rightarrow\!\textsf{Room}$,
with a central $\textsf{Node}$ table (the graph vertex) that a $\textsf{Room}$
optionally maps to one-to-one, $\textsf{Edge}$ rows forming the navigation
graph, and $\textsf{Panorama}$ rows attaching tour imagery to nodes. The
tour's scene list, the minimap's coordinates, and the pathfinder's graph
are three views of this same data; there is no separate graph file. This
shared model is what makes the integration (objective~O and constraint~P5)
concrete rather than nominal: a node added for wayfinding is immediately a
candidate tour scene and a minimap point.

\subsection{Key Data Contracts}
The chat contract is
$\{message,\ session\_id?,\ language?\}\rightarrow\{answer,\ status,\
confidence,\ confidence\_level,\ selected\_agent,\ tool\_used,\
sources[],\ navigation?,\ panorama?,\ session\_id\}$.
The navigation contract is
$\{start\_node\_id,\ destination\_node\_id,\ accessible\_only\}\rightarrow
\{path\_node\_ids,\ path\_node\_names,\ total\_distance,\
estimated\_walk\_time\_minutes,\ is\_accessible,\ turn\_by\_turn[]\}$.
Every AI-generated answer carries its grounding sources as a structural
field, and those sources are always built from retrieved-chunk metadata,
never parsed from the model's text.
